% D.83 candidate. UNREGISTERED, NOTES ONLY, pending full independent audits.
\subsection{A local image criterion with zero compact kernel}

\begin{statement}{D.83}{Open}{Candidate: one local generator for the entire away-five lattice}
Let \(s=(a,b)\) be any completed Kummer pair over \(\mathbb Q\),
with arbitrary infinite ramification, and write its prime coefficients as
\[
 \mathbf v_q=(a_q,b_q)\in\mathbb Z_5^2,
 \qquad \mathbf r=\mathbf v_5.
\]
Define the closed away-five valuation lattice and its local unit vector by
\[
 \begin{split}
 L&=\overline{\sum_{q\ne5}\mathbb Z_5\mathbf v_q}
                                      \ \subset\mathbb Z_5^2,\\
 \lambda_q&=\frac{\log(q^4)}{4\log6}\in\mathbb Z_5,
 \qquad
 \mathbf z=\sum_{q\ne5}\lambda_q\mathbf v_q\in L.
 \end{split}                                                   \tag{D.83.1}
\]
The sums converge in the usual five-adic topology. Suppose
\[
                           L=\mathbb Z_5\mathbf z.              \tag{D.83.2}
\]
The zero lattice is allowed; a nonzero lattice need not be
saturated in \(\mathbb Z_5^2\). For the actual Alexander module
\(A_s\) of D.35 and every finite coefficient quotient
\(W_j=A_s/(5,u,v)^jA_s\), the actual restriction map is an
isomorphism
\[
 H^2(G_{\mathbb Q},W_j)
       \xrightarrow{\ \sim\ }
       \bigoplus_{q\ne5}H^2(G_{\mathbb Q_q},W_j).                 \tag{D.83.3}
\]
Consequently
\[
 \mathcal S_s=\mathcal K_s=0,
 \qquad
 H^2_{\rm cts}(G_{\mathbb Q},A_s)/\overline{\{0\}}
       \xrightarrow{\ \sim\ }\mathcal L_s                       \tag{D.83.4}
\]
topologically, with exactly D.63's local convergence condition.
The actual continuous second cohomology still has D.62's
nonzero nonseparated kernel.

This includes every pair supported only at five and, more generally,
\[
 (a,b)=\mathbf u\,c+\mathbf r\,\kappa(5),                         \tag{D.83.5}
\]
where \(\mathbf u,\mathbf r\in\mathbb Z_5^2\) are arbitrary,
\(c\) has zero five-valuation, its away-five coefficients generate
\(\mathbb Z_5\), and its local unit logarithm coefficient is a unit.
No independence condition on \(\mathbf u,\mathbf r\) is required.
The actual rational point \(t=5\) has this property:
\[
                  (a,b)=(\kappa(5),2\kappa(2)).                  \tag{D.83.6}
\]
Its metabelian fibre is nonempty and is a torsor under
\(\prod_{\mathbb N}\mathbb Z_5\), by D.74.
There is no assertion of section existence for all pairs in
(D.83.2), no general infinitely ramified first-cohomology rank,
and no explicit first higher-derived obstruction.
\end{statement}

\begin{proof}
\paragraph{The infinite sums and their exact cochain representatives.}
The pro-five completion of \(\mathbb Q^\times\) has no sign factor.
In the prime coordinates used in D.35 it is
\(c_0(\{q:q\text{ prime}\},\mathbb Z_5)\). In particular, for
each \(n\), only finitely many vectors \(\mathbf v_q\) fail to
lie in \(5^n\mathbb Z_5^2\). Choose once and for all compatible
roots for each rational prime, giving integral continuous
Kummer cocycles \(\kappa(q)\). They all take values in
\(\mathbb Z_5(1)\), so the sums
\[
 \mathbf s(g)=\mathbf r\kappa(5)(g)
           +\sum_{q\ne5}\mathbf v_q\kappa(q)(g)                 \tag{D.83.7}
\]
converge uniformly in \(g\). They are therefore actual continuous
cocycles, representing the prescribed classes. This argument
does not replace the infinitely ramified class by a finite sum.

Every \(q\ne5\) satisfies \(q^4\in1+5\mathbb Z_5\), and
\(v_5(\log6)=1\). Hence all \(\lambda_q\) in (D.83.1)
are integral and uniformly bounded. The sum defining
\(\mathbf z\) converges and belongs to the closed module \(L\).
The latter is a finitely generated \(\mathbb Z_5\)-module,
although infinitely many primes may contribute to it.

Locally at five, the completed unit class of \(q\) is
\(\lambda_q\kappa(6)\). Thus there is an integral
\(h_q\in\mathbb Z_5\) for which the difference of these two
actual local cocycles is the coboundary of \(h_q\).
The elements \(h_q\) are automatically uniformly bounded.
Therefore
\[
                       \mathbf h=\sum_{q\ne5}\mathbf v_qh_q
                                  \ \in L                       \tag{D.83.8}
\]
converges. A single integral monomial gauge by \(\mathbf h\)
makes the local equality
\[
              \mathbf s(g)=\mathbf z\kappa(6)(g)
                                +\mathbf r\kappa(5)(g)          \tag{D.83.9}
\]
hold for actual cocycles. The gauge lies in \(L\), a point
needed to preserve the global operator subgroup below.

\paragraph{The actual global image and its graph complement.}
Let \(I_0\) be the compact actual untwisted operator image from
D.78. D.82 proves that the actual Kummer cocycle \(\kappa(5)\)
descends continuously to \(I_0\); denote the descended cocycle
by \(k_5\). D.78 also proves that the map
\(G_{\mathbb Q_5}\to I_0\) is surjective and that
\(\kappa(6)\) maps its kernel onto \(\mathbb Z_5\).
These are assertions about the actual arithmetic operators.

The representation of the image of
\[
 g\longmapsto(\mathbf s(g),\rho_0(g))
           \quad\hbox{in }\mathbb Z_5^2\rtimes I_0
\]
by \(t_1^{a(g)}t_2^{b(g)}\rho_0(g)\) is faithful, by D.78:
comparison on \(1,t_1\) determines \(\chi\), the two boundary
axes determine the two translation coordinates, and the
remaining operator determines the untwisted component.
Thus its compact image is the actual operator group \(I_s\).

The untwisted action is unramified at every \(q\ne5\).
Five-primary tame inertia at such a prime has translation
image \(\mathbb Z_5\mathbf v_q\); the chosen generator may
multiply the displayed vector by a unit, which does not
change that image. The closed global image consequently
contains the entire pure translation subgroup \(L\).
Conversely (D.83.7) shows that every translation coordinate,
minus \(\mathbf r k_5\) of its untwisted component, belongs
to \(L\). Surjectivity onto \(I_0\), followed by removal of
this coordinate using the contained subgroup \(L\), proves
the exact compact subgroup identity
\[
 I_s=\{(\mathbf l+\mathbf r k_5(g),g):
                       \mathbf l\in L,\ g\in I_0\}.             \tag{D.83.10}
\]
The set \(\{(\mathbf r k_5(g),g)\}\) is an actual subgroup
because \(k_5\) is a cyclotomic cocycle. It is the graph
complement; it has not been replaced by zero translations.
Conjugation on \(L\) is multiplication by \(\chi\).
A monomial gauge in \(L\), such as (D.83.8), preserves
the subgroup in (D.83.10).

Equation (D.83.9), the local surjection onto \(I_0\),
and the exact image of \(\kappa(6)\) on its kernel now
give the equally exact local identity
\[
 I_5=\{(h\mathbf z+\mathbf r k_5(g),g):
                            h\in\mathbb Z_5,\ g\in I_0\}.      \tag{D.83.11}
\]
Hypothesis (D.83.2) makes (D.83.10) and (D.83.11) equal.
This argument includes nonprimitive lattices and the zero
lattice; no division by a generator of \(L\) occurs.

\paragraph{One decomposition image detects every finite global class.}
We record the finite-level consequence explicitly. Let \(M\)
be any finite continuous module whose global action factors
through \(I_s\), and let \(J\) be its actual finite action
image. A class in
\[
 \operatorname{Sha}^1(\mathbb Q,M)
    =\ker\{H^1(G_{\mathbb Q},M)
                         \longrightarrow\prod_vH^1(G_{\mathbb Q_v},M)\}
\]
has an actual cocycle representative \(z\). This cocycle
vanishes on the kernel of \(G_{\mathbb Q}\to J\).
Indeed, if \(z(g)=m\ne0\) for an element acting trivially
on \(M\), its finite affine image contains \((m,1)\).
The exact finite-image Chebotarev theorem supplies an
unramified prime whose Frobenius is conjugate to that element
\cite[V, Theorem 3.23]{PBC_MilneCFT}. Its action on \(M\)
is the identity, and its cocycle value is nonzero: conjugation
replaces \(m\) by its image under the invertible linear
action of an element of \(J\).
A local coboundary on that Frobenius is zero, a contradiction.
Thus \(z\) descends to a cocycle on \(J\).

Equality \(I_5=I_s\) says that \(G_{\mathbb Q_5}\to J\)
is surjective. Inflation
\(H^1(J,M)\to H^1(G_{\mathbb Q_5},M)\) is injective:
an inflated cocycle which is the coboundary of \(m\)
is already that coboundary on every element of \(J\).
Local triviality at five therefore makes the descended
class zero. We have proved at its original finite level
\[
                       \operatorname{Sha}^1(\mathbb Q,M)=0.
                                                                  \tag{D.83.12}
\]
Also \(M^{G_{\mathbb Q}}=M^{G_{\mathbb Q_5}}\), directly
from equality of their action images.

Apply these facts to the finite Cartier dual \(M=W_j^D\).
Its action factors through \(I_s\), which retains \(\chi\).
Finite global duality identifies
\(\operatorname{Sha}^2(\mathbb Q,W_j)\) with the dual of
\(\operatorname{Sha}^1(\mathbb Q,W_j^D)\), hence it is zero
\cite[I, Theorem 4.10(a)]{DE_MilneADT}.
Finite local duality and equality of the invariant groups
make the component at five
\[
 H^2(G_{\mathbb Q_5},W_j)
            \longrightarrow H^0(G_{\mathbb Q},W_j^D)^*
\]
an isomorphism. The two exact rows of D.63.13 consequently
have zero kernel \(K_j\), proving (D.83.3). In particular
this proves a finite-level isomorphism before taking
any inverse limit. D.63's topological inverse-limit
comparison then gives (D.83.4).

\paragraph{The same conclusion in boundary-depth coordinates.}
For completeness, D.77's explicit three-axis map gives
the same compact vanishing. Let \(e:\mathbb Z_5^2\to\mathbb Z_5\)
be any of the three forms \((x,y)\mapsto x,y,x-y\).
The ideal of global prime coefficients of \(e(s)\) is
generated by \(e(\mathbf r)\) and \(e(L)\). Its local
ideal at five is generated by \(e(\mathbf r)\) and
\(e(\mathbf z)\). They are equal under (D.83.2).
Thus the global and local divisibility exponents agree
on each nonzero axis; on a zero axis the two maps are
the same identity prescribed in D.77. No free compact
factor has been lost in (D.83.4).

\paragraph{Explicit finite and infinite ramification.}
For (D.83.5), write \(c=\sum_{q\ne5}c_q\kappa(q)\).
The condition on its coefficients gives
\(L=\mathbb Z_5\mathbf u\). If its local coefficient
\(\lambda(c)=\sum\lambda_qc_q\) is a unit, then
\(\mathbf z=\lambda(c)\mathbf u\) generates this lattice.
This includes \(\mathbf u=0\) and arbitrary nonprimitive
\(\mathbf u\). It also gives infinitely ramified examples:
for any distinct primes \(q_i\notin\{2,5\}\), take
\[
                       c=\kappa(2)+\sum_{i\ge1}5^i\kappa(q_i).
                                                                  \tag{D.83.13}
\]
The series converges in the completed Kummer space and
\(\lambda(c)\equiv\lambda_2\equiv2\pmod5\), by D.78's
exact logarithm calculation. Its away-five coefficients
generate \(\mathbb Z_5\), as required.

For \(t=5\), the Kummer coordinates are
\(\kappa(5),\kappa(-4)=2\kappa(2)\); the sign has zero
pro-five class. Here \(L=\mathbb Z_5(0,2)\) and
\(\mathbf z=(0,2\lambda_2)\) generates it. The full point
section supplies actual existence, and D.74 supplies
the entire nonempty metabelian fibre. Neither equality
of the operator images nor the separated vanishing is
used to assert arbitrary-pair lifting: the continuous
nonseparated obstruction is still present in general.
\end{proof}
